\section{Introduction}
\label{sec:introduction}
Let $S=\langle m=g_1<g_2<\cdots<g_e=M\rangle$ be a numerical
semigroup, written with its minimal generators. Over a field $k$, its
completed semigroup ring has tangent cone
\[
 G_S=\gr_{\mathfrak m} k[[t^S]],\qquad
 \mathfrak m=(t^{g_1},\ldots,t^{g_e}).
\]
The initial forms of the minimal generators give a surjection
$P_S=k[X_1,\ldots,X_e]\twoheadrightarrow G_S$.
We write
\[
 \beta_i(G_S;k)=\dim_k\Tor_i^{P_S}(G_S,k)
\]
for its total Betti numbers. Thus $\beta_0=1$, and the positive rows
measure the numbers of generators and higher syzygies of the defining
ideal, with a shift of one in the homological index.

The interval completion of $S$ is
\[
 \widetilde S=\langle m,m+1,\ldots,M\rangle.
\]
Its listed generators need not be minimal. Throughout, Betti numbers
on the right are computed over the polynomial ring on the minimal
generators of $\widetilde S$, independently of the presentation of $S$.

\begin{theorem}\label{thm:main}
For every numerical semigroup $S$, every field $k$, and every $i\ge0$,
\[
 \beta_i(G_S;k)\le\beta_i(G_{\widetilde S};k).
\]
\end{theorem}

No restriction is imposed on the characteristic of $k$ or on the
Cohen--Macaulayness of $G_S$. The use of different minimal polynomial
rings is essential: the interval completion may have many more minimal
generators than the original semigroup. The theorem concerns the total
Betti numbers; it does not assert a comparison in each internal degree.

Theorem~\ref{thm:main} resolves Problem~14 in the problem list of
Moscariello--Sammartano \cite{ProblemList}. Herzog--Stamate had
formulated the first-row comparison and discussed its higher-syzygy
extension \cite{APCone}; the problem list states the all-row form used
here.

The target has an explicit Betti table. Put $w=M-m$. If $m\ge2$, its
minimal generators are $m,m+1,\ldots,m+v$, where
$v=\min\{w,m-1\}$. The arithmetic-sequence formula of
Gimenez--Sengupta--Srinivasan, together with the tangent-cone comparison
of Herzog--Stamate, determines its Betti numbers
\cite{APFormula,APCone}. We use that formula as a numerical capacity
against which to compare the source.

The comparison also yields a sharp bound in terms of width alone,
settling Conjecture~1.3 of Caviglia--Moscariello--Sammartano
\cite{CMSBounds}. For the complete local ring
$R_S=k[[t^S]]$, let $\beta_i(R_S;k)$ denote its Betti numbers over
the minimal regular presentation $k[[X_1,\ldots,X_e]]$.
\begin{corollary}\label{cor:width-bound}
For every numerical semigroup $S$ of width $w$, every field $k$, and
every $i\ge1$,
\[
 \beta_i(R_S;k)\le\beta_i(G_S;k)
                 \le i\binom{w+1}{i+1}.
\]
For each $w\ge1$ these bounds are attained simultaneously by
$S=\langle w+1,w+2,\ldots,2w+1\rangle$.
\end{corollary}
The proof, including the local-ring comparison, is given at the end of
Section~\ref{sec:preliminaries}.

The proof separates at $w=m-1$. For $w\ge m-1$, the completion is
the ordinary semigroup $\{0\}\cup[m,\infty)$, and its $i$th positive
Betti number is $i\binom{m}{i+1}$. An Ap\'ery quotient of length $m$
organizes both the regular part and the torsion of $G_S$. One sequence
of linear cuts yields a single resource profile controlling all rows.
The estimates are numerical after that construction; low cut ranks and
large cut ranks require different inequalities.

For $w\le m-2$, divide the offsets $g_j-m$ by their common divisor
$d$. This gives a normalized width $u=w/d$, but leaves the original
semigroup and its physical width $du$ unchanged. Keeping the two endpoint
generators separates the proof into a finite internal core and an
endpoint kernel. The kernel is determined by all competing
representations, and its socle supplies the correction to the core Betti
bound. Estimates for finite cyclic groups and lex quotients then close
the wide normalized ranges. The remaining cases involve one or two
internal variables and are handled directly.

Lan--Tu--Vu proved the interval-completion comparison for three
generators \cite{LanTuVu}. The proof below includes a compatible
one-internal-variable argument, so that the small-generator boundary
uses the same mechanism as the general narrow case. Our use of Green's
restriction theorem, lex Betti bounds, Kneser's theorem, and the
Savchev--Chen theorem is stated where each enters the argument.

The accompanying Lean development verifies the theorem for the actual
completed tangent cones and their minimal presentations. The mathematical
proof is given in Sections~\ref{sec:preliminaries}--\ref{sec:synthesis};
Appendix~\ref{sec:formalization} identifies the formal statement and the
source version used here.

\subsection*{AI Declaration}
Generative AI was used substantially in the mathematical research and
preparation of this manuscript. OpenAI ChatGPT 6 Astra Pro was used to
explore and construct candidate proofs. OpenAI GPT-6, with the
\emph{extra high} reasoning setting, was used to audit the proposed
arguments, look for errors, and check exceptional cases. OpenAI ChatGPT
Work, using GPT-6.1 Sol with the \emph{Ultra} reasoning setting, was
used for later-stage mathematical proof construction, refinement,
integration, and audit. OpenAI Codex
with GPT-6.1 Sol, using the \emph{high} and \emph{Ultra} reasoning
settings, was used to refine and simplify the proofs, develop the Lean
formalization, draft the English manuscript, and revise its exposition.
The model and setting names identify the interfaces used during the work.

The author's contribution was supervisory: he selected the problem,
designed and supervised an iterative workflow of AI proof construction
and separate AI audit, and directed the organization, simplification,
and presentation of the resulting arguments. The AI contributions
included mathematical reasoning and proof development as well as
writing. Audit rounds within this workflow are AI-assisted checks,
rather than independent human referee reports. The author is responsible
for the content and accuracy of the manuscript. The role of Lean and
the verification code is described separately in
Appendix~\ref{sec:formalization}.
